Since this thesis uses promise theory to explain several components of Kubernetes and the architecture of the developed solution, this chapter serves as a brief introduction to the concepts of promise theory.
\cite{promise-theory}, titled \enquote{Promise Theory} is a book by Mark Burgess and Jan A.
Bergstra which introduces the concept of promise theory.

This thesis uses the terms, definitions and notation that can be found in \cite{promise-theory}.
The following notation denotes a promise of the first kind, namely a direct promise from $A$ to $B$ with body $b$ \parencite{promise-theory}:

\[ A \xrightarrow{b} B \]

The body $b$ may contain a label $\Lambda (b)$ uniquely identifying the promise, a type $\tau (b)$ explaining the promise and a constraint $\chi (b)$ on the state of the agent.
Additionally, the body $b$ may be signed to indicate whether the promise is a promise to give or a promise to accept \parencite{promise-theory}:

\[ A \xrightarrow{-b} B \]

Therefore, the promise above is the promise $\Lambda (b)$ from $A$ to $B$ to accept $\tau (b)$ subject to $\chi (b)$.
